\paragraph{短期实验目标}通过"二分法"不断迭代寻找合适的参数

\section{可控参数}

电压、喷射距离、滚筒转速、挤出速率、滑台移速、nZr：PVP配比

需要注意的是，电压和喷射距离本质上是服务于同一个参数: 胶液飞行速度。这两个参数应当整体调节。

\paragraph{预想机理}通过控制胶液飞行速率避免胶液在飞行途中聚团，从而减少附着纤维中的聚团现象；挤出速率：避免在飞行速率过慢时挤出较多导致胶液束过粗；滑台移动速率：减少滑台变速抖动，抑制胶液偏聚；配比：控制纤维的分布；滚筒转速：尽快分离前后附着胶液的空间分布。

\section{实验条件改良}

铝箔纸换硅油纸、提高箱体气密性

\section{实验计划}

\begin{tabular}{|ccc|}
    \hline
    次序&参数&最初预设调参节点\\\hline
    1&nZr: PVP&1.5、2.5\\\hline
    2&滚筒转速&200+n40rpm\\\hline
    3&滑台移速&10mm/h\\\hline
\end{tabular}

电压与喷射距离（喷射速率）的调节需要进行数学计算

\input{maintext/一阶段实验文件/0321改良效果.tex}
\input{maintext/一阶段实验文件/0401效果.tex}
\input{maintext/一阶段实验文件/0522xrd.tex}
\input{maintext/一阶段实验文件/0610实验数据.tex}

\section{材料性能对比}
\input{maintext/文献综述/EAB.tex}
